% Generated draft: requires mathematical and source review.
\documentclass[11pt]{article}
\usepackage[margin=25mm]{geometry}
\usepackage{fontspec}
\setmainfont{Latin Modern Roman}
\usepackage{amsmath,amssymb,mathtools,unicode-math}
\setmathfont{Latin Modern Math}
\usepackage{enumitem,needspace,xurl,hyperref}
\setlength{\parindent}{0pt}
\setlength{\parskip}{5pt}
\setlength{\emergencystretch}{3em}
\newcommand{\N}{\mathbb{N}}
\newcommand{\Z}{\mathbb{Z}}
\newcommand{\Q}{\mathbb{Q}}
\newcommand{\R}{\mathbb{R}}
\newcommand{\C}{\mathbb{C}}
\newcommand{\F}{\mathbb{F}}
\newcommand{\A}{\mathbb{A}}
\newcommand{\PP}{\mathbb{P}}
\newcommand{\E}{\mathbb{E}}
\newcommand{\eps}{\varepsilon}
\newcommand{\dd}{\,\mathrm{d}}
\begin{document}
\section*{UnsolvedMath ID: 30005350}
\textbf{Source problem code:} catalogue entry
\subsection{Definitions and mathematical statement}
The supplied clean statement concerns graph discrepancy, but the supplied original statement concerns balanced slices of a Boolean cube. These formulations are incompatible. The graph formulation below is therefore conditional on explicitly stated conventions; its correspondence with the cited report requires human verification.

For this conditional formulation, let \(G=(V,E)\) be a finite simple undirected graph with \(n=|V|\geq 2\). Define its density and positive discrepancy by
\[
p=\frac{|E|}{\binom{n}{2}},\qquad
\operatorname{disc}^{+}(G)=\max_{U\subseteq V}\left(e_G(U)-p\binom{|U|}{2}\right),
\]
where \(e_G(U)\) is the number of edges with both endpoints in \(U\). These definitions are conventional interpretations, not definitions supplied in the record. Interpret the exceptional densities as
\[
D=\left\{\frac{i-1}{i}:i\in\Z,\ i\geq 1\right\}.
\]
Under these conventions, the proposed statement is
\[
\begin{gathered}
\forall\varepsilon>0\ \exists a=a(\varepsilon)>0\ \forall n\in\Z,\ n\geq2\ \forall G\text{ as above},\\
\left(\forall i\in\Z,\ i\geq1:\ \left|p-\frac{i-1}{i}\right|\geq\varepsilon\right)
\ \Longrightarrow\
\operatorname{disc}^{+}(G)\geq a\sqrt{p(1-p)}\,n^{3/2}.
\end{gathered}
\]
Thus \(a\) is independent of \(n\), \(p\), and \(G\). The separation condition includes equality at distance \(\varepsilon\); the proposed lower bound is also non-strict. The supplied statement gives no sufficiently-large-\(n\) qualification. The graph class, normalization of density and discrepancy, and range of \(i\) must be checked before this conditional formulation is accepted.

The incompatible original passage contains the labels ``Conjecture 8'' and ``Question 7'' and the name Ehud Friedgut. It describes a sub-balanced slice of dimension \(2k\) by fixing \(n-k\) coordinates to \(1\), fixing \(n-k\) coordinates to \(0\), and allowing all assignments to the remaining \(2k\) coordinates having exactly \(k\) ones and \(k\) zeros. It then asks whether every indicator function \(f\) on ``the slice'' with \(\mathbb E f=1/2\) has a sub-balanced slice of dimension \(\log^{*}n\), or of growing dimension, on which \(\mathbb E f>0.51\). The ambient slice and probability measures are not explicitly defined in the passage. Moreover, its dimensions are defined as \(2k\), whereas the requested dimension \(\log^{*}n\) has no rounding or parity convention; ``or growing'' supplies no precise growth quantifiers.

As motivation, that passage says that the parity set has half of every subcube inside it and half outside it, while modulo \(3\) conditions work only on sufficiently large subcubes, and asks for examples on a slice. These are supplied source claims with unspecified scope, not hypotheses of the graph statement or independently verified assertions. The record does not establish how this passage, its attribution, or its numbering relates to the graph conjecture.
\subsection{Short English statement}
The proposed graph conjecture asks for a uniform positive-discrepancy lower bound when the edge density stays at least \(\varepsilon\) away from every density \((i-1)/i\). The displayed formulation uses explicit provisional conventions. The supplied original passage instead asks a question about Boolean slices, so the intended source statement requires human review.
\subsection{Sources}
\begin{itemize}
\item Supplied local UnsolvedMath registry record, permanent ID \(30005350\), title ``Positive Discrepancy away from Multipartite Densities,'' exact source problem code \texttt{catalogue entry}. The proposed year is supplied as 2023, and the graph record has no named proposer.
\item Supplied report citation: \emph{Combinatorics} (2023), Oberwolfach Reports, DOI 10.4171/owr/2023/1, pp.\ 5--89. Supplied source URL: \url{https://doi.org/10.4171/owr/2023/1}. The record describes Oberwolfach Reports as documenting mathematical workshops and their open-problem discussions. No problem-specific page is supplied.
\item Supplied literature source, labeled ``Oberwolfach Reports 2023/1'': \url{https://doi.org/10.4171/owr/2023/1}. Literature assessment dated exactly \texttt{2026-08-22}: the supplied curation labels the graph conjecture open and reports finding no primary result proving or disproving the stated density-separated positive-discrepancy lower bound. This is supplied dated curation, not an independently verified literature assessment; the record itself requests independent verification of time-sensitive status and cited claims.
\item The supplied record labels its graph statement \texttt{corrected\_verified} and says it was checked against the official report and preserves its hypotheses, notation, and target. This is a supplied verification claim, not verification performed for this draft. It conflicts with the unrelated Boolean-slice text in the original-statement field.
\item The original-statement field contains the name Ehud Friedgut adjacent to problem labels and a Boolean-slice question. This does not establish attribution of the graph conjecture to him.
\end{itemize}
\end{document}
